\documentclass[11pt]{article}
\usepackage[margin=1in]{geometry}
\usepackage{amsmath,amssymb,amsthm}
\usepackage{booktabs}
\usepackage{hyperref}

\newtheorem{lemma}{Lemma}
\newtheorem{proposition}{Proposition}
\theoremstyle{definition}
\newtheorem{definition}{Definition}

\newcommand{\conj}{\mathrel{\sim}}
\newcommand{\id}{\mathrm{id}}

\title{Disproof of Conjecture 00000005400:\\
Orbit Distributions, Period Counts, and the Impossibility of\\
a Conjugate Pair with Different Periods}
\author{TLMC submission --- verdict: \textbf{FALSE}}
\date{September 13, 2026}

\begin{document}
\maketitle

\begin{abstract}
Conjecture 00000005400 asserts two things at once: (i) there exist two maps with
identical orbit distribution limits but different periodic point counts, and
(ii) the separation is realized by an explicit \emph{conjugate pair} with the
same measure but different periods. We show that both layers fail in the literal
discrete reading. First, on finite sets the periodic point count is a function of
the orbit-length distribution, $P(f)=\sum_{l\ge1} l\,N_l(f)$, so layer (i) is
self-contradictory. Second, exact conjugacy preserves pointwise periods,
$\operatorname{per}_{\tau f\tau^{-1}}(\tau x)=\operatorname{per}_f(x)$, so the
witness objects demanded by layer (ii) --- a conjugate pair with different
periods --- \emph{do not exist}. We verify both identities by exhaustive machine
enumeration: all $2304$ triples $(\tau,f,x)\in S_4\times S_4\times\{0,1,2,3\}$
for conjugacy, and all $4^4=256$ self-maps of a $4$-point set for the counting
identity. A Lean~4 development with zero axioms and zero \texttt{sorry}
re-derives everything by kernel-checked computation.
\end{abstract}

\section{The conjecture}

\begin{quote}
\emph{Definition: Orbit distribution and period counting are two layers.
Conjecture: There exist two maps with identical orbit distribution limits but
different periodic point counts, and the separation is realized by an explicit
conjugate pair with the same measure but different periods.}
\end{quote}

The conjecture treats ``orbit distribution'' and ``period counting'' as two
independent layers that can come apart, and claims the coming-apart is exhibited
by a conjugate pair. We refute both layers, in the natural discrete reading in
which a \emph{map} is a self-map of a set, the orbit distribution records cycle
lengths, and conjugacy means $g=\tau f\tau^{-1}$ for a bijection $\tau$.

\section{Definitions}

Let $X$ be a set and $f\colon X\to X$.

\begin{definition}[orbit, periodic point, period]\label{def:period}
For $x\in X$ write $f^0x=x$ and $f^{k+1}x=f(f^kx)$. The point $x$ is
\emph{periodic} if $f^kx=x$ for some $k\ge 1$; the least such $k$ is the
\emph{period} $\operatorname{per}_f(x)$. The \emph{periodic orbit} of a periodic
point $x$ is $\mathcal{O}_f(x)=\{f^kx : k\ge 0\}$; all its elements are periodic
with the same period, and $|\mathcal{O}_f(x)|=\operatorname{per}_f(x)$.
\end{definition}

\begin{definition}[orbit distribution, periodic point count]\label{def:profile}
For $l\ge 1$ let $N_l(f)$ be the number of distinct periodic orbits of $f$ of
length exactly $l$, and let $P(f)$ be the number of periodic points of $f$. The
\emph{orbit distribution} of $f$ (its orbit profile) is the multiset
$\{1^{N_1(f)},2^{N_2(f)},3^{N_3(f)},\dots\}$, equivalently the sorted list of
cycle lengths.
\end{definition}

\begin{definition}[conjugacy]\label{def:conj}
Two maps $f,g\colon X\to X$ are \emph{conjugate}, written $f\conj g$, if
$g=\tau f\tau^{-1}$ for some bijection $\tau\colon X\to X$.
\end{definition}

\section{Layer 1 fails: the orbit distribution determines the periodic point count}

\begin{proposition}[Identity I]\label{prop:I}
For every self-map $f$ of a finite set,
\begin{equation}
P(f)\;=\;\sum_{l\ge 1} l\,N_l(f).
\end{equation}
Consequently, if $f$ and $g$ have the same orbit distribution then
$P(f)=P(g)$. There are no two maps with identical orbit distribution but
different periodic point counts.
\end{proposition}

\begin{proof}
By Definition~\ref{def:period} the periodic points of $f$ are partitioned into
the distinct periodic orbits $\mathcal{O}$ of $f$, and each orbit of length $l$
contains exactly $l$ points, each of period $l$. Counting periodic points orbit
by orbit, each orbit of length $l$ is counted exactly once as an orbit in $N_l$
and contributes exactly $l$ points:
\[
P(f)=\Bigl|\bigsqcup_{\mathcal{O}\text{ orbit of }f}\mathcal{O}\Bigr|
=\sum_{\mathcal{O}}|\mathcal{O}|=\sum_{l\ge1}\sum_{\substack{\mathcal{O}\\
|\mathcal{O}|=l}}l=\sum_{l\ge1}l\,N_l(f).
\]
The right-hand side is a function of the orbit distribution alone, so equal
distributions force equal counts.
\end{proof}

Identity I makes the first half of the conjecture self-contradictory: the two
``layers'' are not independent; the second is a function of the first.

\paragraph{Computational check.}
Over all $4^4=256$ self-maps of $X=\{0,1,2,3\}$, the identity
$P(f)=\sum l\,N_l(f)$ holds in all $256$ cases, and among the $11$ distinct
orbit profiles no profile is realized with two different values of $P$
(Table~\ref{tab:selfmaps}). The script \texttt{reproduce.py} (part~B) re-runs
this audit.

\begin{table}[ht]
\centering
\begin{tabular}{lrr}
\toprule
orbit profile & periodic count $P$ & number of maps\\
\midrule
$(1)$          & 1 & 64\\
$(1,1)$        & 2 & 48\\
$(1,1,1)$      & 3 & 12\\
$(1,1,1,1)$    & 4 & 1\\
$(1,1,2)$      & 4 & 6\\
$(1,2)$        & 3 & 36\\
$(1,3)$        & 4 & 8\\
$(2)$          & 2 & 48\\
$(2,2)$        & 4 & 3\\
$(3)$          & 3 & 24\\
$(4)$          & 4 & 6\\
\bottomrule
\end{tabular}
\caption{All $256$ self-maps of $\{0,1,2,3\}$ grouped by orbit profile. Each of
the $11$ profiles carries exactly one periodic point count: the profile
determines $P(f)=\sum l\,N_l(f)$.}
\label{tab:selfmaps}
\end{table}

\section{Layer 2 fails: conjugacy preserves periods, so the demanded witnesses do not exist}

\begin{proposition}[Identity II]\label{prop:II}
Let $\tau\colon X\to X$ be a bijection and $g=\tau f\tau^{-1}$. Then for every
$x\in X$ and every $k\ge0$,
$g^k(\tau x)=\tau(f^kx)$. In particular
\begin{equation}
\operatorname{per}_g(\tau x)=\operatorname{per}_f(x)\quad\text{for all }x\in X.
\end{equation}
So period is a conjugacy invariant, and no conjugate pair has different periods.
\end{proposition}

\begin{proof}
Induction on $k$. For $k=0$ both sides equal $\tau x$. If
$g^k(\tau x)=\tau(f^kx)$, then
$g^{k+1}(\tau x)=g\bigl(g^k(\tau x)\bigr)=\tau f\tau^{-1}\bigl(\tau(f^kx)\bigr)
=\tau\bigl(f(f^kx)\bigr)=\tau(f^{k+1}x)$,
because $\tau^{-1}\tau=\id$. Now $g^k(\tau x)=\tau x$ iff $\tau(f^kx)=\tau x$,
and since $\tau$ is injective this holds iff $f^kx=x$. Taking the least such
$k\ge1$ on both sides gives $\operatorname{per}_g(\tau x)=\operatorname{per}_f(x)$.
\end{proof}

Proposition~\ref{prop:II} uses nothing about finiteness: for \emph{any} two
conjugate maps, the number of points of each exact period agrees. Therefore the
second half of the conjecture --- ``the separation is realized by an explicit
conjugate pair \dots\ with different periods'' --- demands objects whose defining
description is unsatisfiable. Whenever the first half could conceivably hold
(non-conjugate maps sharing an orbit-distribution limit), the pair realizing it
is, by that very clause, required to be conjugate --- and conjugate pairs cannot
separate. The conjunction is false under every reading.

\paragraph{Computational check.}
On $X=\{0,1,2,3\}$ we enumerated all triples
$(\tau,f,x)\in S_4\times S_4\times\{0,1,2,3\}$: all $2304$ satisfy
$\operatorname{per}_{\tau f\tau^{-1}}(\tau x)=\operatorname{per}_f(x)$, zero
violations. As an illustration, take $f=(0\,1)(2\,3)$, $\tau=(1\,2\,3)$,
$g=\tau f\tau^{-1}=(0\,2)(1\,3)$; then every point of $f$ and of $g$ has period
$2$, and the pointwise comparison agrees on all four points
(Table~\ref{tab:pair}).

\begin{table}[ht]
\centering
\begin{tabular}{cccc}
\toprule
$x$ & $\operatorname{per}_f(x)$ & $\tau(x)$ &
$\operatorname{per}_{\tau f\tau^{-1}}(\tau x)$\\
\midrule
0 & 2 & 0 & 2\\
1 & 2 & 2 & 2\\
2 & 2 & 3 & 2\\
3 & 2 & 1 & 2\\
\bottomrule
\end{tabular}
\caption{A concrete conjugate pair $f=(0\,1)(2\,3)$, $g=\tau f\tau^{-1}$ with
$\tau=(1\,2\,3)$: pointwise periods are transported exactly.}
\label{tab:pair}
\end{table}

For reference, the five cycle types of $S_4$ (partitions of $4$) all give
$P=4$, as every point of a permutation is periodic:

\begin{table}[ht]
\centering
\begin{tabular}{lrrrrr}
\toprule
cycle type (profile) & $(1,1,1,1)$ & $(1,1,2)$ & $(1,3)$ & $(2,2)$ & $(4)$\\
\midrule
number of permutations & 1 & 6 & 8 & 3 & 6\\
periodic points $P$    & 4 & 4 & 4 & 4 & 4\\
\bottomrule
\end{tabular}
\caption{Cycle types of $S_4$: $24$ permutations, $5$ profiles, one periodic
count each.}
\label{tab:s4}
\end{table}

\section{Machine verification}

The companion \texttt{lean4/} project (Lean $4.33.1$, no Mathlib, no classical
axioms) formalizes the $4$-point case with permutation-valued and
map-valued encodings:

\begin{itemize}
\item \texttt{conj\_preserves\_period} : Identity II, proved by \texttt{decide}
over $24\times256\times4=24576$ configurations ($\tau$ over all permutations,
$f$ over \emph{all} self-maps, $x$ over all points);
\item \texttt{count\_from\_profile\_sum} : Identity I's bridge
$P(f)=(\text{profile}(f)).\mathrm{sum}$, decided over all $256$ self-maps;
\item \texttt{profile\_det\_count} : equal profiles force equal periodic counts,
proved by rewriting with the bridge lemma;
\item \texttt{permAtTable\_eq} : kernel-checked certificate that the indexed
maps \texttt{permAt} enumerate exactly the $24$ permutations;
\item \texttt{conjWitnessExists\_false} and \texttt{refute} : the complete
search for a conjugate pair with a pointwise period difference evaluates to
\texttt{false}, and the conjecture's existential is negated by pure logic from
\texttt{conj\_preserves\_period}.
\end{itemize}

\texttt{Check.lean} prints \texttt{does not depend on any axioms} for every
theorem: the development is axiom-free in the strongest sense --- not even
\texttt{propext} or \texttt{Quot.sound} is used --- and \texttt{sorry}-free,
with all enumeration performed by the kernel.

\section{Conclusion}

\begin{enumerate}
\item[(a)] The periodic point count is determined by the orbit distribution
(Proposition~\ref{prop:I}); the conjecture's first layer cannot separate.
\item[(b)] Period is a conjugacy invariant (Proposition~\ref{prop:II}); the
conjecture's second layer demands conjugate witnesses with different periods,
which do not exist.
\end{enumerate}

The conjecture 00000005400 is therefore \textbf{false}: the alleged two-layer
separation is internally inconsistent, and its realization by a conjugate pair is
logically impossible.

\end{document}
